\documentclass[12pt, a4paper]{article}

% === Paquetes ===
\usepackage{fontspec}
\usepackage{polyglossia}
\setmainlanguage{spanish}
\usepackage{geometry}
\usepackage{booktabs}
\usepackage{enumitem}
\usepackage{hyperref}
\usepackage{xcolor}
\usepackage{amsmath}
\usepackage{fancyhdr}
\usepackage{titlesec}
\usepackage{caption}
\usepackage{float}
\usepackage{listings}
\usepackage{tikz}
\usetikzlibrary{shapes.geometric, arrows.meta, positioning, fit, backgrounds}

% === Configuración ===
\geometry{margin=2.5cm}
\hypersetup{
    colorlinks=true,
    linkcolor=black!70!blue,
    urlcolor=black!70!blue,
    citecolor=black!70!blue
}

\definecolor{accent}{RGB}{30, 90, 160}
\definecolor{codebg}{RGB}{248, 248, 248}
\definecolor{codegreen}{RGB}{40, 120, 60}

\titleformat{\section}
  {\Large\bfseries\color{accent}}{\thesection}{1em}{}
\titleformat{\subsection}
  {\large\bfseries\color{accent!80!black}}{\thesubsection}{1em}{}

\pagestyle{fancy}
\fancyhf{}
\fancyhead[L]{\small\textcolor{accent}{InvestigIA -- Diseño de integración}}
\fancyhead[R]{\small\textcolor{accent}{\thepage}}
\renewcommand{\headrulewidth}{0.4pt}

\lstdefinestyle{python}{
    language=Python,
    backgroundcolor=\color{codebg},
    basicstyle=\ttfamily\small,
    keywordstyle=\color{accent}\bfseries,
    commentstyle=\color{codegreen},
    stringstyle=\color{red!60!black},
    breaklines=true,
    frame=single,
    framerule=0.3pt,
    rulecolor=\color{black!30},
    xleftmargin=0.5cm,
    numbers=left,
    numberstyle=\tiny\color{black!40},
    numbersep=8pt,
    showstringspaces=false,
    tabsize=4
}

\lstdefinestyle{json}{
    backgroundcolor=\color{codebg},
    basicstyle=\ttfamily\small,
    breaklines=true,
    frame=single,
    framerule=0.3pt,
    rulecolor=\color{black!30},
    xleftmargin=0.5cm,
    showstringspaces=false
}

% === Documento ===
\begin{document}

% --- Título ---
\begin{center}
    {\LARGE\bfseries Diseño de integración}\\[0.3em]
    {\LARGE\bfseries de Gentle-AI en InvestigIA}\\[1.5em]
    {\normalsize
    Adaptación de Engram, Skills, Contratos y Flujo SDD}\\
    {\normalsize al ecosistema multiagente de InvestigIA}\\[1em]
    Centro de Investigación en Computación, IPN\\[0.5em]
    {\normalsize \today}
\end{center}

\vspace{2em}

% ======================================================================
\section{Introducción}
% ======================================================================

InvestigIA transforma un perfil de estudiante de posgrado en un protocolo
preliminar de investigación mediante seis agentes especializados: perfilado,
matching, metodológico, investigador, redacción y revisión. El notebook del
proyecto establece que cada agente debe cumplir con contrato de datos, prompt
versionado, máximo de iteraciones, manejo de errores, logging y pruebas.

El problema es que esos requisitos quedan enunciados pero no diseñados. ¿Qué
pasa exactamente cuando el agente de perfil genera un JSON incompleto? ¿Cómo
sabe el agente de redacción qué evidencia recuperó el investigador? ¿Con qué
criterio se decide si el protocolo final es suficiente o necesita otra ronda
de correcciones?

Este documento propone una respuesta concreta tomando los conceptos de
Gentle-AI (memoria persistente, flujo estructurado con validación, instrucciones
especializadas por contexto) y adaptándolos a la arquitectura que ya define
InvestigIA. No se trata de instalar Gentle-AI como dependencia externa, sino
de implementar las mismas ideas directamente en el código del proyecto.

% ======================================================================
\section{Flujo multiagente con puertas de validación}
% ======================================================================

El pipeline de InvestigIA no debería ser lineal. Si el agente de perfil
genera un JSON con la lista de intereses vacía, no tiene sentido pasar eso
al agente de matching. Lo razonable es validar la salida de cada agente
antes de permitirle avanzar al siguiente.

A esas validaciones les llamamos puertas. Cada puerta $V_i$ revisa que la
salida de un agente cumpla un contrato mínimo antes de pasarla al siguiente
nodo del flujo.

\begin{figure}[H]
\centering
\resizebox{\textwidth}{!}{
\begin{tikzpicture}[
    node distance=0.9cm and 1cm,
    agent/.style={
        rectangle, rounded corners=3pt, minimum width=2cm, minimum height=0.85cm,
        draw=accent, thick, fill=accent!8, font=\footnotesize\bfseries, text=accent
    },
    gate/.style={
        diamond, aspect=1.8, minimum width=0.9cm, minimum height=0.7cm,
        draw=orange!70!black, thick, fill=orange!10, font=\scriptsize\bfseries, text=orange!70!black
    },
    input/.style={
        ellipse, minimum width=1.6cm, minimum height=0.7cm,
        draw=gray, thick, fill=gray!10, font=\footnotesize
    },
    outcome/.style={
        rectangle, rounded corners=2pt, minimum width=1.5cm, minimum height=0.6cm,
        draw=codegreen!60!black, thick, fill=codegreen!10, font=\footnotesize, text=codegreen!40!black
    },
    arrow/.style={-{Stealth[length=4pt]}, thick, draw=black!60},
    failarrow/.style={-{Stealth[length=4pt]}, thick, draw=red!50, dashed},
]

% Fila 1: izquierda a derecha
\node[input] (est) {Estudiante};
\node[agent, right=0.5cm of est] (perf) {Perfil};
\node[gate, right=0.4cm of perf] (v1) {$V_1$};
\node[agent, right=0.4cm of v1] (match) {Matching};
\node[gate, right=0.4cm of match] (v2) {$V_2$};
\node[agent, right=0.4cm of v2] (met) {Metodológico};

% Fila 2: derecha a izquierda (serpentina)
\node[agent, below=1.4cm of met] (inv) {Investigador};
\node[gate, left=0.4cm of inv] (v3) {$V_3$};
\node[agent, left=0.4cm of v3] (red) {Redacción};
\node[gate, left=0.4cm of red] (v4) {$V_4$};
\node[agent, left=0.4cm of v4] (rev) {Revisión};
\node[gate, left=0.4cm of rev] (v5) {$V_5$};
\node[outcome, left=0.5cm of v5] (ok) {Entrega};

% Flechas principales
\draw[arrow] (est) -- (perf);
\draw[arrow] (perf) -- (v1);
\draw[arrow] (v1) -- (match);
\draw[arrow] (match) -- (v2);
\draw[arrow] (v2) -- (met);
\draw[arrow] (met) -- (inv);
\draw[arrow] (inv) -- (v3);
\draw[arrow] (v3) -- (red);
\draw[arrow] (red) -- (v4);
\draw[arrow] (v4) -- (rev);
\draw[arrow] (rev) -- (v5);
\draw[arrow] (v5) -- (ok);

% Flechas de reintento (fila 1)
\draw[failarrow] (v1.south) -- ++(0,-0.35) -| ([xshift=-0.2cm]perf.south);
\node[font=\scriptsize, text=red!60!black, below=0.2cm of v1, anchor=north] {reintenta};

\draw[failarrow] (v2.south) -- ++(0,-0.35) -- ++(-0.8,0) |- ([xshift=-0.2cm]match.south);
\node[font=\scriptsize, text=red!60!black, below=0.2cm of v2, anchor=north] {reintenta};

% Flechas de reintento (fila 2)
\draw[failarrow] (v3.north) -- ++(0,0.35) -| ([xshift=0.2cm]inv.north);
\node[font=\scriptsize, text=red!60!black, above=0.2cm of v3, anchor=south] {reintenta};

\draw[failarrow] (v4.north) -- ++(0,0.35) -- ++(0.8,0) |- ([xshift=0.2cm]red.north);
\node[font=\scriptsize, text=red!60!black, above=0.2cm of v4, anchor=south] {reintenta};

\draw[failarrow] (v5.south) -- ++(0,-0.35) -| ([xshift=-0.2cm]rev.south);
\node[font=\scriptsize, text=red!60!black, below=0.2cm of v5, anchor=north] {itera};

% Etiquetas de las puertas
\node[font=\scriptsize, text=black!50, below=0.05cm of v1] {contrato};
\node[font=\scriptsize, text=black!50, below=0.05cm of v2] {contrato};
\node[font=\scriptsize, text=black!50, above=0.05cm of v3] {metadatos};
\node[font=\scriptsize, text=black!50, above=0.05cm of v4] {evidencia};
\node[font=\scriptsize, text=black!50, below=0.05cm of v5] {coherencia};

\end{tikzpicture}
}
\caption{Pipeline multiagente con puertas de validación y caminos de reintento}
\label{fig:pipeline}
\end{figure}

Las reglas de cada puerta:

\begin{description}[style=nextline, leftmargin=2.5cm]
    \item[$V_1$: Perfil $\to$ Matching]
    El perfil JSON debe tener nombre no vacío, nivel académico válido
    (licenciatura, maestria o doctorado), y al menos un interés. Si falta
    cualquiera de estos campos, el agente de perfil vuelve a preguntarle
    al estudiante.

    \item[$V_2$: Matching $\to$ Metodológico]
    El ranking debe contener al menos un resultado con score mayor a cero.
    Un matching vacío significa que los intereses del estudiante no
    coinciden con ninguna línea o asesor disponible; en ese caso se le
    pide que amplíe o reformule.

    \item[$V_3$: Metodológico $\to$ Investigador]
    El esquema de investigación necesita pregunta de investigación, al
    menos un objetivo específico y una estrategia metodológica explícita.
    Si el agente metodológico no logra definir alguno de estos elementos,
    itera con las observaciones que resulten del intento fallido.

    \item[$V_4$: Investigador $\to$ Redacción]
    Cada fragmento recuperado debe incluir identificador, título, autores,
    año y texto. Los fragmentos que no cumplan se descartan. Si después
    del filtrado no queda ninguno, la evidencia se marca como
    insuficiente y el agente de redacción lo recibe como input explícito.

    \item[$V_5$: Revisión $\to$ Entrega]
    El agente de revisión evalúa el protocolo contra reglas deterministas
    y genera un puntaje de coherencia entre 0 y 1. Si el puntaje cae por
    debajo del umbral (0.7 en la primera implementación), el protocolo
    regresa a redacción junto con las observaciones específicas. Si pasa,
    se entrega al estudiante.
\end{description}

La idea de las puertas viene del concepto de verificación en SDD: cada fase
debe confirmarse antes de avanzar. En un sistema multiagente esto es
particularmente útil porque un error en un nodo temprano se propaga a todos
los posteriores.

% ======================================================================
\section{Engram como capa de estado persistente}
% ======================================================================

Gentle-AI usa Engram para guardar decisiones de arquitectura entre sesiones
de coding. En InvestigIA el concepto se adapta para algo distinto: mantener
el estado completo de la interacción con el estudiante.

La diferencia es sutil pero importante. En Gentle-AI, Engram guarda
decisiones que tomó el desarrollador. En InvestigIA, Engram guarda el
resultado del trabajo que produjeron los agentes para un estudiante
específico.

\newpage
\subsection{Qué se persiste}

\begin{table}[h!]
\centering
\begin{tabular}{@{}p{3.5cm}p{4cm}p{5.5cm}@{}}
\toprule
\textbf{Dato} & \textbf{Topic key} & \textbf{Quién lo escribe} \\
\midrule
Perfil del estudiante & \texttt{session/{id}/perfil} & Agente de Perfil \\
\addlinespace
Líneas seleccionadas & \texttt{session/{id}/matching} & Agente de Matching \\
\addlinespace
Esquema de investigación & \texttt{session/{id}/esquema} & Agente Metodológico \\
\addlinespace
Evidencia recuperada & \texttt{session/{id}/evidencia} & Agente Investigador \\
\addlinespace
Protocolo generado & \texttt{session/{id}/protocolo} & Agente de Redacción \\
\addlinespace
Reporte de revisión & \texttt{session/{id}/revision} & Agente de Revisión \\
\addlinespace
Historial de iteraciones & \texttt{session/{id}/iteraciones} & Orquestador \\
\bottomrule
\end{tabular}
\caption{Estructura de persistencia en Engram}
\end{table}

\subsection{Para qué sirve}

Recuperación de sesión. Si el estudiante cierra la aplicación a la mitad
del proceso y regresa horas después, el sistema retoma exactamente donde
quedó. El perfil, la evidencia ya recuperada, el esquema metodológico,
todo sigue ahí. Sin esta persistencia, cada sesión empieza de cero.

Contexto entre agentes. Cuando el agente de redacción necesita generar el
protocolo, consulta directamente en Engram la evidencia recuperada y el
esquema metodológico. El orquestador no tiene que pasar toda esa
información como parámetro en cada llamada. Los agentes leen lo que
necesitan de la capa de estado compartido.

Auditoría. Cada iteración queda registrada con timestamp y resultado.
Si el agente de revisión rechaza el protocolo tres veces seguidas, el
historial permite analizar qué cambió entre cada intento. ¿Mejoró la
redacción? ¿Se agregaron fuentes? ¿O el agente simplemente repite los
mismos errores?

Métricas. El conteo de iteraciones por agente responde directamente a la
pregunta de investigación PI-5. Si el agente metodológico necesita en
promedio 2.3 iteraciones para producir un esquema válido, mientras que
el de perfil pasa en 1.1, eso indica dónde el sistema tiene más
dificultades y qué componentes aportan menos valor del esperado.

\newpage
\subsection{Modelo de datos}

Cada observación en Engram sigue esta estructura:

\begin{lstlisting}[style=json]
{
  "topic_key": "session/abc123/perfil",
  "type": "state",
  "content": {
    "nombre": "María López",
    "nivel": "maestria",
    "intereses": ["robótica", "visión por computadora"],
    "experiencia": "proyecto de detección de objetos",
    "timestamp": "2026-06-17T10:30:00Z"
  },
  "session_id": "abc123"
}
\end{lstlisting}

% ======================================================================
\section{Skills por agente}
% ======================================================================

En Gentle-AI, un skill es un conjunto de instrucciones que el agente de
coding carga cuando el contexto de trabajo lo requiere. En InvestigIA se
aplica la misma idea, pero las instrucciones definen el comportamiento de
cada agente del sistema.

Los skills no son prompts. El prompt es lo que se envía al LLM en cada
llamada. El skill es la capa encima: las reglas, validaciones y patrones
que el agente debe seguir independientemente del prompt específico que
reciba en una iteración determinada.

\subsection{Agente de Perfil}

\begin{description}[style=nextline, leftmargin=3cm]
    \item[\texttt{perfil/estructura}]
    Define el esquema JSON del perfil con sus campos obligatorios y
    opcionales. La validación se hace con Pydantic: intereses mínimo 1
    elemento, nombre no vacío, nivel restringido a tres valores posibles.

    \item[\texttt{perfil/preguntas-adaptativas}]
    Lógica para generar preguntas de seguimiento cuando el perfil está
    incompleto. Si el estudiante no menciona experiencia previa, el
    agente pregunta directamente por eso. Si los intereses son demasiado
    amplios (``inteligencia artificial'' sin más), pide que delimite.
\end{description}

\subsection{Agente de Matching}

\begin{description}[style=nextline, leftmargin=3cm]
    \item[\texttt{matching/scoring}]
    La función de puntuación del notebook: $S(a_i, s) = w_1 I_i + w_2
    E_i + w_3 H_i + w_4 D_i$. Los pesos $w_j$ son configurables y deben
    sumar 1. En la PoC se puede empezar con pesos iguales (0.25 cada uno)
    y ajustar según los resultados de los casos de prueba.

    \item[\texttt{matching/explainability}]
    Cada resultado del ranking incluye una explicación textual. El
    estudiante necesita saber \emph{por qué} un asesor es candidato, no
    solo que aparece primero en la lista. Sin explicación, el matching es
    una caja negra.
\end{description}

\subsection{Agente Metodológico}

\begin{description}[style=nextline, leftmargin=3cm]
    \item[\texttt{metodologico/delimitacion}]
    Flujo guiado para construir el esquema de investigación: problema,
    alcance, pregunta, objetivos (general y específicos), hipótesis.
    Cada sección tiene sus propias reglas de validación interna.

    \item[\texttt{metodologico/coherencia-interna}]
    Reglas deterministas que se evalúan antes de pasar al siguiente
    agente. La pregunta de investigación debe conectarse con el objetivo
    general. Los objetivos específicos deben ser alcanzables con la
    metodología propuesta. La hipótesis, si existe, debe ser testeable.
\end{description}

\subsection{Agente Investigador (RAG)}

\begin{description}[style=nextline, leftmargin=3cm]
    \item[\texttt{investigador/retrieval}]
    Pipeline de recuperación: a partir del esquema de investigación se
    genera una consulta expandida, se calcula similitud coseno sobre los
    embeddings del corpus, y se recuperan los 5 fragmentos más relevantes.
    Cada fragmento conserva sus metadatos completos.

    \item[\texttt{investigador/metadatos}]
    Validación estricta de metadatos. Identificador, título, autores,
    año y texto son obligatorios. Un fragmento sin autores o con año
    vacío se descarta sin excepción. Si todos se descartan, la evidencia
    se marca como insuficiente.

    \item[\texttt{investigador/citas}]
    El agente no puede generar citas que no existan en los fragmentos
    recuperados. Esta es probablemente la regla más importante del
    sistema: si el LLM necesita apoyar una afirmación y no hay evidencia
    disponible, debe marcarlo explícitamente como \texttt{[sin fuente]}
    en lugar de inventar una referencia.
\end{description}

\subsection{Agente de Redacción}

\begin{description}[style=nextline, leftmargin=3cm]
    \item[\texttt{redaccion/protocolo}]
    Plantilla del protocolo con sus secciones mínimas: pregunta,
    objetivos, hipótesis, metodología, justificación, alcance,
    resultados esperados. Cada sección se genera condicionada por la
    evidencia que recuperó el agente investigador.

    \item[\texttt{redaccion/separacion-evidencia}]
    Cada afirmación del protocolo debe distinguir entre lo que proviene
    de la evidencia recuperada y lo que el LLM generó por su cuenta. Las
    afirmaciones sin fuente se marcan como \texttt{[sin fuente]}. El
    estudiante y el asesor deben poder diferenciar a simple vista qué
    tiene respaldo documental y qué no.
\end{description}

\subsection{Agente de Revisión}

\begin{description}[style=nextline, leftmargin=3cm]
    \item[\texttt{revision/reglas-deterministas}]
    Lista de verificación automática: ¿la pregunta conecta con el
    objetivo general? ¿Los objetivos específicos son medibles? ¿La
    metodología es compatible con el tipo de estudio? ¿Las referencias
    citadas existen en el corpus recuperado?

    \item[\texttt{revision/puntaje}]
    Genera un puntaje de coherencia entre 0 y 1. El umbral de aprobación
    inicial es 0.7, aunque habría que calibrarlo con los casos de prueba
    C1 a C6 del diseño experimental.

    \item[\texttt{revision/iteracion}]
    Máximo 3 intentos. Si después de tres revisiones el puntaje sigue
    por debajo del umbral, se entrega el protocolo con las observaciones
    pendientes y una nota indicando que requiere revisión humana. Es
    preferible entregar algo imperfecto con las fallas documentadas que
    entrar en un bucle infinito.
\end{description}

% ======================================================================
\section{Contratos entre agentes}
% ======================================================================

Cada transición entre agentes tiene un contrato definido con Pydantic. La
idea es simple: si la salida de un agente no cumple el contrato que espera
el siguiente, el orquestador captura el error y itera en lugar de propagar
datos inválidos por todo el pipeline.

\subsection{Contrato del agente de perfil}

\begin{lstlisting}[style=python]
from pydantic import BaseModel, Field
from typing import List
from enum import Enum

class NivelAcademico(str, Enum):
    licenciatura = "licenciatura"
    maestria = "maestria"
    doctorado = "doctorado"

class PerfilEstudiante(BaseModel):
    nombre: str = Field(..., min_length=1)
    nivel: NivelAcademico
    intereses: List[str] = Field(..., min_items=1)
    experiencia_previa: str | None = None
    institucion: str | None = None
\end{lstlisting}

Cuando el LLM genera un JSON con \texttt{intereses: []}, Pydantic levanta
un error de validación. El orquestador captura ese error, le pregunta al
estudiante qué líneas de investigación le interesan, y vuelve a invocar
el agente de perfil con la respuesta.

\subsection{Contrato del agente investigador}

\begin{lstlisting}[style=python]
from pydantic import BaseModel, Field

class FragmentoRecuperado(BaseModel):
    id: str
    titulo: str = Field(..., min_length=1)
    autores: List[str] = Field(..., min_items=1)
    anio: int = Field(..., ge=1900, le=2030)
    texto: str = Field(..., min_length=50)
    score: float = Field(..., ge=0.0, le=1.0)

class EvidenciaRecuperada(BaseModel):
    fragmentos: List[FragmentoRecuperado]
    evidencia_insuficiente: bool = False

    def tiene_evidencia_valida(self) -> bool:
        return len(self.fragmentos) > 0
\end{lstlisting}

Los fragmentos sin autores o con año fuera de rango se descartan en la
validación. Si no queda ninguno, \texttt{evidencia\_insuficiente} se marca
como \texttt{True}. El agente de redacción recibe esa señal y, en lugar
de inventar referencias, genera el protocolo indicando qué secciones
carecen de respaldo documental.

\subsection{Contrato de revisión}

\begin{lstlisting}[style=python]
from pydantic import BaseModel, Field
from typing import List

class Observacion(BaseModel):
    seccion: str
    tipo: str  # "error" | "warning" | "sugerencia"
    descripcion: str

class ReporteRevision(BaseModel):
    puntaje_coherencia: float = Field(..., ge=0.0, le=1.0)
    observaciones: List[Observacion]
    aprobado: bool

    @property
    def necesita_iteracion(self) -> bool:
        return not self.aprobado
\end{lstlisting}

El orquestador lee \texttt{necesita\_iteracion}. Cuando es \texttt{True},
pasa las observaciones al agente de redacción como contexto adicional.
El agente de redacción entonces sabe qué secciones fallaron y por qué,
lo cual es más útil que simplemente volver a intentar sin información
sobre el error anterior.

% ======================================================================
\section{Orquestación con LangGraph}
% ======================================================================

LangGraph modela el flujo como un grafo dirigido donde cada nodo es un
agente y cada arista es una puerta de validación. Las aristas condicionales
permiten que el grafo decida en tiempo de ejecución si avanza o itera.

\begin{lstlisting}[style=python]
from langgraph.graph import StateGraph

# Estado compartido entre todos los nodos
class EstadoInvestigIA(TypedDict):
    perfil: PerfilEstudiante | None
    matching: MatchingResultado | None
    esquema: EsquemaInvestigacion | None
    evidencia: EvidenciaRecuperada | None
    protocolo: Protocolo | None
    revision: ReporteRevision | None
    iteraciones: dict  # conteo por agente

# Grafo
grafo = StateGraph(EstadoInvestigIA)

grafo.add_node("perfil", agente_perfil)
grafo.add_node("matching", agente_matching)
grafo.add_node("metodologico", agente_metodologico)
grafo.add_node("investigador", agente_investigador)
grafo.add_node("redaccion", agente_redaccion)
grafo.add_node("revision", agente_revision)

# Aristas condicionales: validación
grafo.add_conditional_edges(
    "perfil",
    validar_perfil,        # ¿cumple contrato?
    {"ok": "matching", "retry": "perfil"}
)
grafo.add_conditional_edges(
    "revision",
    validar_protocolo,     # ¿puntaje >= 0.7?
    {"ok": END, "retry": "redaccion"}
)
\end{lstlisting}

Cada función de validación (como \texttt{validar\_perfil} o
\texttt{validar\_protocolo}) evalúa el contrato Pydantic correspondiente
y devuelve ``ok'' o ``retry''. El grafo no necesita lógica condicional
escrita a mano; LangGraph se encarga del enrutamiento.

\subsection{Límite de iteraciones}

El notebook establece ``máximo de iteraciones'' como requisito de cada
agente. La implementación propone 3 intentos como límite. Cuando un agente
falla tres veces, el sistema registra las tres salidas y sus errores en
Engram, marca esa etapa como \texttt{requiere\_intervencion\_humana}, y
continúa con el siguiente agente si el flujo lo permite.

Esto evita que el sistema se quede atascado en un bucle sin salida. En un
sprint de 10 días, es prefericiente detectar dónde el LLM no puede solo
y documentar la limitación, en vez de gastar tiempo ajustando prompts
para un caso que tal vez requiera intervención humana de todas formas.

% ======================================================================
\section{Integración con LM Studio}
% ======================================================================

Los seis agentes se comunican con Gemma 4 a través de LM Studio usando
la interfaz compatible con OpenAI. El cliente se instancia una vez y se
comparte:

\begin{lstlisting}[style=python]
from openai import OpenAI
import os

cliente = OpenAI(
    base_url=os.getenv("LLM_BASE_URL", "http://localhost:1234/v1"),
    api_key=os.getenv("LLM_API_KEY", "lm-studio"),
)

def invocar_agente(
    sistema: str,
    usuario: str,
    modelo: str = "gemma-4",
    temperatura: float = 0.2,
) -> str:
    respuesta = cliente.chat.completions.create(
        model=modelo,
        temperature=temperatura,
        messages=[
            {"role": "system", "content": sistema},
            {"role": "user", "content": usuario},
        ],
    )
    return respuesta.choices[0].message.content
\end{lstlisting}

Cada skill define su propio prompt de sistema. La temperatura varía
según la tarea: 0.1 o 0.2 para perfil y revisión, donde se busca
determinismo; 0.4 o 0.6 para redacción, donde cierta variabilidad
mejora la calidad del texto generado.

% ======================================================================
\section{Mapeo con las preguntas de investigación}
% ======================================================================

El diseño propuesto no es independiente de las preguntas de investigación
del notebook. Cada componente responde a una o más preguntas de forma
directa.

\begin{table}[h!]
\centering
\begin{tabular}{@{}p{1.5cm}p{5cm}p{5.5cm}@{}}
\toprule
\textbf{PI} & \textbf{Pregunta} & \textbf{Cómo se aborda} \\
\midrule
PI-1 & Multiagente transforma información no estructurada en protocolo & Flujo con 6 agentes y contratos Pydantic \\
\addlinespace
PI-2 & Trazabilidad bibliográfica con RAG local & Fragmentos con metadatos obligatorios, citas solo del corpus \\
\addlinespace
PI-3 & Latencias compatibles con interacción & Gemma 4 para todo excepto verificación \\
\addlinespace
PI-4 & Validaciones deterministas + LLM & Puertas $V_i$ con reglas Pydantic + revisión por LLM \\
\addlinespace
PI-5 & Qué componentes aportan más valor & Engram registra iteraciones por agente, se puede medir \\
\bottomrule
\end{tabular}
\caption{Relación entre preguntas de investigación y diseño}
\end{table}

La pregunta PI-3 sobre latencias es la que menos se aborda en este
documento. La asignación de Gemma 4 como modelo principal ayuda, pero
las latencias reales dependen del hardware, la cuantización y el tamaño
del corpus. Eso habría que medirlo en la PoC, no se puede resolver en
el diseño.

% ======================================================================
\section{Limitaciones}
% ======================================================================

Gemma 4 y validación estricta. Los contratos Pydantic asumen que el LLM
genera JSON válido. Un modelo cuantizado puede producir JSON malformado,
especialmente en respuestas largas. La mitigación es combinar reintentos
con un prompt de corrección y un parser tolerante que intente reparar
errores menores (comas finales, comillas simples en vez de dobles).

Dependencia de Engram. Si Engram no está disponible en el entorno de
despliegue, el sistema necesita un fallback. La opción más simple es
almacenamiento local: un archivo JSON por sesión en disco, o una tabla
SQLite. Funciona, pero pierde la ventaja de la memoria compartida entre
sesiones.

Umbral de revisión. El puntaje 0.7 para aprobar el protocolo es una
primera aproximación. Podría ser demasiado permisivo o demasiado
estricto; no hay forma de saberlo sin ejecutar los casos de prueba C1
a C6 y ver qué puntajes producen. Es un hiperparámetro que se calibra
en la práctica.

Costo de iteración. Cada intento adicional es una llamada extra a Gemma 4.
Si un agente falla tres veces, eso son tres llamadas que no estaban en
el flujo ideal. En el hardware institucional, cada llamada puede tomar
entre 10 y 30 segundos dependiendo de la longitud del prompt. No es
crítico, pero sí se nota en un sprint de 10 días.

Skills como código. Los skills descritos aquí son descripciones
conceptuales. Para que el orquestador los cargue dinámicamente, hay que
implementarlos como archivos con un formato estándar (el mismo que usa
Gentle-AI para sus skills). Esa implementación no está hecha todavía.

% ======================================================================
\section{Próximos pasos}
% ======================================================================

\begin{enumerate}[leftmargin=1.5em]
    \item Definir los esquemas Pydantic completos para los 6 agentes.
    \item Implementar el grafo LangGraph con las aristas condicionales.
    \item Escribir los prompts de sistema de cada agente (primer versión).
    \item Configurar Engram con la estructura de topic keys propuesta.
    \item Probar con el caso C1 (estudiante con tema definido) del
    diseño experimental.
    \item Calibrar el umbral de revisión con los casos C1 a C6.
\end{enumerate}

% ======================================================================
% Referencias
% ======================================================================
\newpage
\begin{thebibliography}{9}

\bibitem{investigia-notebook}
Equipo InvestigIA,
\textit{Plan de acción para el desarrollo de InvestigIA},
Notebook Jupyter, CIC-IPN, 2026.

\bibitem{gentle-ai}
Gentleman Programming,
\textit{Gentle-AI -- Ecosystem, Frameworks, Workflows for AI coding agents},
\url{https://github.com/Gentleman-Programming/gentle-ai},
2026.

\bibitem{engram}
Gentleman Programming,
\textit{Engram Commands -- Gentle-AI},
\url{https://github.com/Gentleman-Programming/gentle-ai/blob/main/docs/engram.md},
2026.

\bibitem{langgraph}
LangChain,
\textit{LangGraph -- Building agent state machines},
\url{https://github.com/langchain-ai/langgraph},
2025.

\end{thebibliography}

\end{document}
